% 对应题号 M2-027，本题条件、变量及几何关系逐项核对。
% 修订：启用中文断行；核对题设及图文对应。
\documentclass[tikz,border=14pt]{standalone}
\usepackage{fontspec,xeCJK}
\usepackage{amsmath,amssymb}
\usepackage{xcolor}
\usepackage{tikz}
\usetikzlibrary{arrows.meta,calc,patterns,positioning,bending}
\IfFontExistsTF{Heiti SC}{\IfFontExistsTF{Heiti SC}{\setmainfont{Heiti SC}}{\setmainfont{FandolHei-Regular}}\IfFontExistsTF{Heiti SC}{\setsansfont{Heiti SC}}{\setsansfont{TeX Gyre Heros}}}{\IfFontExistsTF{Songti SC}{\setmainfont{Songti SC}}{\setmainfont{FandolSong-Regular}}\IfFontExistsTF{Songti SC}{\setsansfont{Songti SC}}{\setsansfont{FandolSong-Regular}}}
\definecolor{bluefield}{RGB}{24,76,255}
\definecolor{redfield}{RGB}{232,49,35}
\definecolor{orangefield}{RGB}{236,132,18}
\definecolor{greenfield}{RGB}{18,132,86}
\definecolor{purplefield}{RGB}{112,72,190}
\definecolor{linegray}{RGB}{60,68,78}
\definecolor{softblue}{RGB}{238,243,255}
\definecolor{softgreen}{RGB}{238,249,244}
\definecolor{softorange}{RGB}{255,247,235}
\tikzset{
  arrow/.style={-{Latex[length=3.1mm,width=2.1mm]},line width=1.1pt},
  thinarr/.style={-{Latex[length=2.5mm,width=1.7mm]},line width=.85pt},
  axis/.style={-{Latex[length=2.5mm,width=1.7mm]},line width=.8pt,draw=linegray},
  guide/.style={draw=linegray!55,dashed,line width=.65pt},
  box/.style={draw=linegray,rounded corners=3pt,fill=white,align=center,inner sep=6pt,line width=.75pt},
  bluebox/.style={box,draw=bluefield,fill=softblue},
  greenbox/.style={box,draw=greenfield,fill=softgreen},
  orangebox/.style={box,draw=orangefield,fill=softorange}
}
\IfFontExistsTF{Heiti SC}{\setCJKmainfont[AutoFakeBold=1.5]{Heiti SC}}{\setCJKmainfont[AutoFakeBold=1.5]{FandolHei-Regular}}
\IfFontExistsTF{Heiti SC}{\setCJKsansfont{Heiti SC}}{\setCJKsansfont{FandolHei-Regular}}
\begin{document}

\begin{tikzpicture}[x=1cm,y=1cm]
\path[use as bounding box] (-10,-4.6) rectangle (10,4.6);
\fill[white] (-10,-4.6) rectangle (10,4.6);
\node[font=\Large\bfseries,text width=18.6cm,align=center] at (0,4.05) {线性感应偶极：极化能含二分之一，粒子趋向强场区};
\begin{scope}[shift={(-4.7,-.1)}]
\draw[axis] (-2.8,-1.6)--(2.9,-1.6) node[right] {$+x$};
\foreach \x/\l in {-2.2/.5,-1.1/.7,0/.9,1.1/1.1}{\draw[arrow,bluefield] (\x,1.3)--({\x+\l},1.3);}
\node[bluefield] at(0,2){$\mathbf E(x)=E_0(1+x/L)\mathbf e_x$};
\fill[redfield] (0,0) circle(.1) node[above left] {粒子};
\draw[arrow,greenfield] (0,0)--(1.7,0) node[midway,below] {$\mathbf p=\alpha\mathbf E$};
\draw[arrow,orangefield] (0,-.95)--(1.8,-.95) node[right] {$\mathbf F$};
\node[orangefield] at (0,-2.3) {$\alpha>0,L>0$：力指向强场侧};
\end{scope}
\begin{scope}[shift={(4.55,-.1)}]
\node[bluebox,text width=6.7cm] at (0,1.65) {$\displaystyle U=-\int_0^E\alpha E'\,dE'=-\frac12\alpha E^2$};
\node[greenbox,text width=6.7cm] at (0,0) {$\displaystyle \mathbf F=-\nabla U=\frac12\alpha\nabla E^2$};
\node[orangebox,text width=6.7cm] at (0,-1.7) {$E(x)=E_0(1+x/L)\Rightarrow F_x(0)=\alpha E_0^2/L$};
\end{scope}
\end{tikzpicture}
\end{document}
